\section{Tổng quan đề tài}

\subsection{Thực trạng chuỗi cung ứng nông sản Việt Nam}

Chuỗi cung ứng nông sản truyền thống tại Việt Nam hiện vẫn phụ thuộc nặng vào các khâu trung gian. Theo khảo sát của Viện Chính sách và Chiến lược Phát triển Nông nghiệp Nông thôn (IPSARD) được dẫn lại trên Tạp chí Công Thương (2023--2024), một loại nông sản trung bình phải trải qua ít nhất 4--5 khâu trung gian trước khi đến tay người tiêu dùng cuối, và sự can thiệp của các khâu này làm tăng giá thành lên 30--40\% so với giá tại trang trại \cite{tapchicongthuong_ipsard_2023_2024}. Cùng nguồn, năm 2023 hệ thống chợ truyền thống vẫn chiếm tới 80\% sản lượng nông sản tiêu thụ tại các thành phố lớn như Hà Nội và TP. Hồ Chí Minh, trong khi hệ thống bán lẻ hiện đại (siêu thị, chuỗi cửa hàng) chỉ chiếm khoảng 15--20\% \cite{tapchicongthuong_ipsard_2023_2024}. Đến năm 2024, tỷ lệ nông sản được sản xuất theo hình thức liên kết chuỗi (giữa nông dân/HTX với doanh nghiệp) mới đạt 21,6\% theo báo cáo của Bộ Nông nghiệp và Phát triển nông thôn \cite{tapchicongthuong_ipsard_2023_2024} --- cho thấy phần lớn sản lượng còn lại vẫn giao dịch qua kênh tự do, nơi thương lái giữ vai trò định giá chi phối.

Vai trò của thương lái trong bối cảnh này mang tính hai mặt: một mặt họ giải quyết bài toán thu gom, sơ chế và vận chuyển phân tán mà từng nông hộ nhỏ khó tự thực hiện; mặt khác, do nắm thông tin thị trường tốt hơn nông dân (bất cân xứng thông tin --- xem mục 3.4) và thường là điểm bán duy nhất tại địa phương vào thời điểm thu hoạch, thương lái có vị thế đàm phán giá vượt trội, dẫn đến hiện tượng ép giá khi vào vụ thu hoạch rộ. Trên một số thị trường xuất khẩu, theo Viện Logistics và Chuỗi cung ứng (VILAS), có tới 85--90\% lượng nông sản Việt Nam vào thị trường thế giới phải đi qua trung gian bằng thương hiệu nước ngoài, và phần giá trị nông dân/doanh nghiệp Việt thực nhận chỉ tương đương khoảng 25\% giá bán đến tay người tiêu dùng cuối \cite{vilas_logistics}.

Nguyên nhân sâu xa của hiện tượng ``được mùa mất giá'' gồm ba nhóm yếu tố chính: (i) cung tăng ồ ạt và đồng loạt khi vào vụ thu hoạch do sản xuất còn theo tập quán mùa vụ, thiếu quy hoạch rải vụ; (ii) thiếu phương án tiêu thụ đa kênh và năng lực chế biến sâu/bảo quản để hấp thụ sản lượng dư thừa; (iii) tỷ lệ hao hụt, hư hỏng sau thu hoạch rất cao --- theo chuyên gia Nguyễn Thanh Bình được TheLeader (2024) dẫn lại, tỷ lệ này với một số loại trái cây có thể lên tới 35--40\% --- buộc nông dân phải bán tháo với giá thấp trước khi hàng hỏng hoàn toàn \cite{theleader_2024}. Đây chính là ba ``điểm đau'' mà một sàn TMĐT như Farmery có thể can thiệp trực tiếp: (i) kết nối cung--cầu theo thời gian thực và mở rộng thị trường tiêu thụ ra ngoài phạm vi địa lý của thương lái địa phương; (ii) cho phép đặt hàng B2B theo lịch/hợp đồng để nông dân chủ động sản xuất; (iii) tối ưu logistics lạnh và rút ngắn thời gian từ thu hoạch đến giao hàng.

Cập nhật số liệu gần nhất (2025--2026): theo phản ánh trên báo Tuổi Trẻ (đầu năm 2026) về vấn đề định giá nông sản, tỷ lệ tổn thất sau thu hoạch của ngành rau quả Việt Nam vẫn ở mức 20--40\%, gây thiệt hại kinh tế khoảng 3,5--4,1 tỷ USD mỗi năm; nguyên nhân trực tiếp là hạ tầng bảo quản lạnh quá mỏng --- cả nước hiện chỉ có khoảng 117 kho lạnh chuyên nghiệp, trong đó khoảng 90\% phục vụ đông lạnh sâu cho thủy sản/thịt, gần như không có kho lạnh chuyên dụng cho rau quả tươi. Do thiếu nơi lưu trữ đệm, nông dân ``buộc phải bán ngay, bán vội'' ngay khi thu hoạch, càng làm suy yếu vị thế đàm phán giá trước thương lái \cite{tuoitre_2026}. Số liệu này củng cố thêm cho luận điểm ở trên và cho thấy vấn đề chưa được cải thiện đáng kể so với năm 2024, đồng thời chỉ ra rằng đầu tư vào logistics lạnh là điều kiện gần như bắt buộc để một sàn TMĐT như Farmery giải quyết triệt để hiện tượng ``được mùa mất giá'' chứ không chỉ dừng ở việc kết nối cung--cầu.

\subsection{Bối cảnh chuyển đổi số nông nghiệp tại Việt Nam}

Chương trình OCOP (Mỗi xã một sản phẩm) là chương trình quốc gia hỗ trợ chuẩn hóa và nâng cao giá trị sản phẩm nông thôn/nông sản đặc trưng vùng miền. Theo số liệu năm 2024 được Dân Việt tổng hợp, cả nước từng có 14.642 sản phẩm đạt chuẩn OCOP từ 3 sao trở lên, tăng 3.586 sản phẩm so với năm 2023; trong đó 73,2\% đạt 3 sao, 23,5\% đạt 4 sao và 51 sản phẩm đạt hạng cao nhất 5 sao, với mục tiêu ban đầu đến hết năm 2025 là chuẩn hóa khoảng 15.000 sản phẩm \cite{danviet_2024_ocop}. Số liệu cập nhật đến hết năm 2025 (VietNamNet, công bố 31/12/2025) cho thấy chương trình đã vượt mục tiêu này: cả nước đạt 18.243 sản phẩm OCOP từ 3 sao trở lên (tăng 13.756 sản phẩm so với năm 2020), trong đó 71,4\% đạt 3 sao, 27,7\% đạt 4 sao và 126 sản phẩm đạt 5 sao; số chủ thể tham gia tăng lên 9.820 đơn vị (33,4\% là HTX, 26,5\% là doanh nghiệp nhỏ, 32,6\% là cơ sở sản xuất/hộ kinh doanh), trong đó khoảng 40\% chủ thể do phụ nữ làm chủ và 17,1\% thuộc đồng bào dân tộc thiểu số \cite{vietnamnet_2025_ocop}. Giai đoạn 2026--2030, chương trình được định hướng chuyển trọng tâm ``từ số lượng sang chất lượng'', trong đó đẩy mạnh số hóa sản xuất, truy xuất nguồn gốc và tiêu thụ qua thương mại điện tử được xác định là một trong các trụ cột \cite{vietnamnet_2025_ocop} --- đây là tín hiệu chính sách rất thuận lợi cho định hướng của Farmery. Đây là nguồn cung sản phẩm có thương hiệu và câu chuyện xuất xứ rất phù hợp để đưa lên một sàn TMĐT có tính năng truy xuất nguồn gốc như Farmery.

Về chính sách, Quyết định số 749/QĐ-TTg ngày 03/6/2020 của Thủ tướng Chính phủ phê duyệt ``Chương trình Chuyển đổi số quốc gia đến năm 2025, định hướng đến năm 2030'' đã xác định nông nghiệp là một trong 8 lĩnh vực cần ưu tiên chuyển đổi số (cùng Y tế, Giáo dục, Tài chính-Ngân hàng, Giao thông vận tải-logistics, Năng lượng, Tài nguyên-Môi trường, Sản xuất công nghiệp), với các mục tiêu cụ thể: phát triển nông nghiệp thông minh và nông nghiệp chính xác; xây dựng hệ thống dữ liệu lớn về đất đai, cây trồng, vật nuôi, thủy sản; tự động hóa quy trình sản xuất và quản lý chuỗi cung ứng; thúc đẩy thương mại điện tử nông sản; và triển khai định hướng ``mỗi nông dân là một thương nhân'' ứng dụng công nghệ số \cite{qd749_2020}. Gần đây hơn, Nghị quyết số 57-NQ/TW của Bộ Chính trị về đột phá phát triển khoa học, công nghệ, đổi mới sáng tạo và chuyển đổi số quốc gia tiếp tục được xác định là động lực cho chuyển đổi số ngành nông nghiệp trong giai đoạn mới \cite{nq57_tw}.

Cập nhật quan trọng về bộ máy và chính sách (2025): từ ngày 01/3/2025, Bộ Nông nghiệp và Phát triển nông thôn đã hợp nhất với Bộ Tài nguyên và Môi trường thành Bộ Nông nghiệp và Môi trường (Bộ NN\&MT) \cite{botnmt_2025}; do đó các báo cáo/quyết định của ngành nông nghiệp ban hành từ thời điểm này mang tên cơ quan mới. Cụ thể, ngày 31/12/2025, Bộ NN\&MT đã ban hành Quyết định số 5943/QĐ-BNNMT phê duyệt kế hoạch chuyển đổi số của ngành, với các mục tiêu: đưa 100\% thủ tục hành chính đủ điều kiện lên trực tuyến toàn trình (tối thiểu 85\% hồ sơ xử lý trực tuyến), người dân/doanh nghiệp chỉ khai báo dữ liệu một lần; xây dựng cơ sở dữ liệu dùng chung, đồng bộ giữa lĩnh vực nông nghiệp và môi trường; và ứng dụng dữ liệu lớn/AI trong giám sát, dự báo sản xuất nông nghiệp \cite{qd5943_2025}. Đây là văn bản chính sách cập nhật nhất, cụ thể hóa thêm cho định hướng đã đặt ra tại Quyết định 749/QĐ-TTg (2020), và là căn cứ chính sách quan trọng, mới nhất để đề tài Farmery khẳng định tính phù hợp với định hướng vĩ mô của Nhà nước.

Về tiêu chuẩn chất lượng, VietGAP (Vietnamese Good Agricultural Practices) là tiêu chuẩn thực hành nông nghiệp tốt do Việt Nam ban hành, gồm khoảng 70 tiêu chí, được công nhận trong nước và phù hợp với các chuỗi cung ứng nội địa như siêu thị, bếp ăn tập thể, chợ đầu mối. GlobalGAP là tiêu chuẩn quốc tế nghiêm ngặt hơn nhiều, với 252 tiêu chuẩn (36 tiêu chí bắt buộc tuân thủ 100\%, 127 tiêu chí yêu cầu tuân thủ tối thiểu 95\%, và 89 tiêu chí khuyến nghị), được công nhận toàn cầu và là điều kiện để nông sản tiếp cận các thị trường khó tính như EU, Nhật Bản, Hàn Quốc, Úc, Mỹ. Sản phẩm đạt GlobalGAP được gán mã GGN 13 chữ số và đăng ký trong cơ sở dữ liệu toàn cầu, giúp truy xuất minh bạch hơn so với VietGAP (chỉ có giấy chứng nhận, chưa có mã định danh chuẩn hóa riêng). Về nguyên tắc, một cơ sở đạt GlobalGAP đã vượt chuẩn VietGAP nên không cần chứng nhận VietGAP riêng, và VietGAP thường được xem là bước đệm năng lực trước khi nông dân/HTX tiến tới GlobalGAP nếu có định hướng xuất khẩu \cite{vietgap_globalgap}. Đây là căn cứ để Farmery thiết kế quy trình kiểm duyệt chứng nhận theo 2 cấp độ (nội địa/VietGAP và xuất khẩu/GlobalGAP) tương ứng với 2 nhóm khách mua lẻ và mua sỉ B2B.

\subsection{So sánh các sàn thương mại điện tử nông sản hiện có}

Tại Việt Nam:

FoodMap là sàn TMĐT nông sản tư nhân hoạt động theo cả mô hình B2C và B2B, hợp tác với hơn 300 nông trại và nhà sản xuất, gắn mã QR trên toàn bộ sản phẩm để khách hàng tra cứu nguồn gốc qua website/ứng dụng. FoodMap đã gọi vốn 500.000 USD (vòng hạt giống, 9/2020) và 2,9 triệu USD (vòng pre-Series A, 12/2021); trong giai đoạn giãn cách vì COVID-19 năm 2021, cả lượng người dùng và doanh số của FoodMap từng tăng đột biến khoảng 300\% \cite{foodmap}.

Postmart.vn (Tổng công ty Bưu điện Việt Nam -- VNPost) và Vỏ Sò (Viettel Post) từng là các sàn TMĐT do doanh nghiệp bưu chính nhà nước vận hành, tập trung hỗ trợ nông dân đưa sản phẩm --- trong đó có nhiều sản phẩm OCOP --- lên môi trường số, tận dụng sẵn hạ tầng logistics bưu chính toàn quốc. Trong một chiến dịch hỗ trợ tiêu thụ nông sản mùa dịch (giai đoạn 2021), hai sàn này đã giúp tiêu thụ gần 1.200 tấn nông sản cho các tỉnh phía Nam \cite{postmart_voso_2021}. Cập nhật quan trọng: cả hai sàn này hiện không còn hoạt động dưới tên cũ. Vỏ Sò (voso.vn) đã ngừng hoạt động từ đầu năm 2023 với lý do ``nâng cấp website'' và chưa vận hành trở lại; Postmart.vn được đổi tên thành Buudien.vn từ ngày 31/3/2024 để định vị thương hiệu chuyên biệt hơn cho nông--lâm--thủy sản \cite{buudien_2024}. Đến tháng 12/2024, cả hai tập đoàn bưu chính-chuyển phát đều ra mắt sàn chuyên biệt mới theo hướng thu hẹp vào B2B/xuất khẩu --- vốn rất gần với định hướng của Farmery: VNPost ra mắt nongsan.buudien.vn, kế thừa Buudien.vn nhưng chuyên biệt cho nông sản chất lượng cao và hỗ trợ nông dân xuất khẩu; Viettel Post ra mắt VIPO Mall, một sàn B2B bán sỉ xuyên biên giới kết nối hai chiều khách hàng Việt Nam với nhà cung cấp quốc tế, tập trung xuất khẩu nông--lâm--thủy sản và hàng thủ công mỹ nghệ \cite{vipo_nongsan_2024}. Điểm mạnh chung của nhóm sàn nhà nước/bưu chính là mạng lưới vận chuyển rộng và được hậu thuẫn truyền thông; điểm hạn chế là các sàn thế hệ trước (Postmart, Vỏ Sò) thiên về hỗ trợ tiêu thụ thời vụ/ứng phó thiên tai-dịch bệnh hơn là nền tảng thương mại có chiến lược sản phẩm dài hạn, còn các sàn thế hệ mới (nongsan.buudien.vn, VIPO Mall) tuy đã hướng tới B2B/xuất khẩu nhưng vẫn thiếu tính năng phục vụ đồng thời cả bán lẻ cá nhân lẫn đàm phán giá sỉ theo MOQ trên cùng một nền tảng như Farmery hướng tới.

Mô hình quốc tế tham khảo:

e-Choupal (Ấn Độ, tập đoàn ITC triển khai từ năm 2000) là mô hình ``click and mortar'' kết nối trực tiếp nông dân với thị trường thông qua các kiosk internet đặt tại làng, loại bỏ tầng trung gian truyền thống (mandi/thương lái địa phương). Theo Business Standard (2020), mô hình này giúp nông dân tăng thu nhập tới 50\%, đã tiếp cận khoảng 4 triệu nông dân, thu mua khoảng 3 triệu tấn nông sản từ 225 huyện thuộc 22 tỉnh, với năng suất tăng thêm 10--25\% nhờ áp dụng khuyến nghị canh tác và ITC tiết kiệm được 6--10\% chi phí thu mua so với kênh truyền thống \cite{business_standard_2020_echoupal}. Đây là một trong những case điển hình nhất trên thế giới về việc số hóa giúp giảm phụ thuộc thương lái mà đề tài Farmery có thể tham khảo về cơ chế định giá minh bạch tại điểm thu mua.

Pinduoduo / Duo Duo Grocery (Trung Quốc) áp dụng mô hình mua chung (group buying) kết hợp bán trực tiếp từ nông dân đến người tiêu dùng, giảm phụ thuộc vào các tầng phân phối trung gian. Theo thông cáo của Nasdaq (2021), tính đến năm 2021 nền tảng này đã kết nối trực tiếp hơn 12 triệu nông dân với người tiêu dùng; riêng năm 2020, tổng giá trị giao dịch (GMV) liên quan nông nghiệp đạt hơn 270 tỷ nhân dân tệ, tăng từ 136 tỷ nhân dân tệ năm 2019 \cite{nasdaq_2021_pinduoduo}. Mô hình ``đặt hàng trước 23h, nhận hàng hôm sau từ nông trại địa phương'' là gợi ý tham khảo tốt cho quy trình vận chuyển hàng tươi của Farmery.

Khoảng trống chung so với đặc thù nông sản: phần lớn các sàn kể trên chỉ mạnh ở một phía --- hoặc thuần B2C (Pinduoduo hướng người tiêu dùng cuối), hoặc thuần thu mua B2B tập trung (e-Choupal chủ yếu phục vụ chuỗi cung ứng nội bộ của ITC; VIPO Mall và nongsan.buudien.vn của Việt Nam cũng đang đi theo hướng chuyên biệt B2B/xuất khẩu) --- mà chưa tích hợp đầy đủ cả hai luồng bán lẻ và bán sỉ theo MOQ/giá bậc trên cùng một nền tảng như định hướng của Farmery. Bên cạnh đó, tính năng truy xuất nguồn gốc bằng mã QR ở các sàn Việt Nam (như FoodMap) hiện chủ yếu phục vụ mục đích minh bạch thông tin cơ bản, chưa gắn liền với nhật ký canh tác chi tiết theo chuẩn quốc tế (GS1 --- xem mục 3.3), và hạ tầng logistics lạnh (cold-chain) cho hàng dễ hư hỏng tại Việt Nam nói chung vẫn còn hạn chế --- như số liệu cập nhật ở mục 1.1 cho thấy, cả nước hiện chỉ có khoảng 117 kho lạnh chuyên nghiệp và gần như không có kho lạnh chuyên dụng cho rau quả tươi \cite{tuoitre_2026}, so với các nền tảng lớn ở Trung Quốc. Đây là những khoảng trống mà Farmery đặt mục tiêu giải quyết: kết hợp truy xuất nguồn gốc gắn nhật ký canh tác, hỗ trợ song song bán lẻ và bán sỉ B2B theo MOQ, và ưu tiên thiết kế quy trình vận chuyển cho hàng tươi dễ hư hỏng.

\subsection{Phân tích 4 nhóm người dùng}

\subsubsection*{a) Nông dân / hợp tác xã (HTX)}
\textbf{Nhu cầu:} bán được sản lượng thu hoạch với giá công bằng hơn kênh thương lái, tiếp cận trực tiếp khách hàng B2B (nhà hàng, siêu thị) để có đơn hàng ổn định theo hợp đồng, xây dựng thương hiệu/uy tín cá nhân qua lịch sử đánh giá và chứng nhận. \textbf{Khó khăn:} hạn chế về kỹ năng số và thiết bị (điện thoại thông minh, kết nối internet ổn định ở vùng sâu vùng xa); thiếu năng lực tự chụp ảnh/ghi nhật ký canh tác đúng chuẩn; tâm lý e ngại công nghệ và lo sợ bị nền tảng ``thu phí thay thương lái'' mà không mang lại giá tốt hơn thực chất.

\subsubsection*{b) Người tiêu dùng cá nhân}
\textbf{Nhu cầu:} mua nông sản tươi, an toàn, có thể tra cứu nguồn gốc rõ ràng, giá cả hợp lý, giao hàng nhanh và đúng chất lượng như hình ảnh/mô tả. \textbf{Khó khăn:} khó phân biệt sản phẩm ``sạch'' thật với sản phẩm gắn mác marketing (do bất cân xứng thông tin --- mục 3.4); lo ngại hàng tươi bị hư hỏng trong vận chuyển; thói quen mua sắm nông sản tại chợ/siêu thị truyền thống vẫn còn phổ biến, đòi hỏi sàn TMĐT phải tạo được lòng tin đủ lớn để thay đổi hành vi.

\subsubsection*{c) Doanh nghiệp thu mua (nhà hàng, siêu thị, nhà phân phối)}
\textbf{Nhu cầu:} nguồn cung ổn định về số lượng và chất lượng theo tiêu chuẩn (VietGAP/GlobalGAP tùy mục đích nội địa hay xuất khẩu), có thể đặt hàng theo hợp đồng/lịch giao định kỳ, cơ chế giá theo sản lượng (giá bậc) và công nợ/thanh toán linh hoạt. \textbf{Khó khăn:} khó tìm và xác minh nhiều nhà cung cấp nhỏ lẻ phân tán; rủi ro đứt gãy cung ứng khi mất mùa hoặc nhà cung cấp không đáp ứng MOQ; thiếu công cụ quản lý hợp đồng và hóa đơn tập trung khi làm việc với nhiều nông dân/HTX riêng lẻ.

\subsubsection*{d) Quản trị viên hệ thống}
\textbf{Nhu cầu:} công cụ kiểm duyệt hiệu quả (nhà bán, sản phẩm, chứng nhận) để đảm bảo chất lượng thông tin trên sàn; công cụ giám sát và xử lý vi phạm (hàng giả nhãn, gian lận chứng nhận, khiếu nại gian dối); báo cáo thống kê để ra quyết định vận hành. \textbf{Khó khăn:} khối lượng kiểm duyệt lớn khi số nhà bán tăng nhanh; khó xác minh tính xác thực của giấy chứng nhận VietGAP/GlobalGAP nếu không có API kết nối với cơ quan cấp chứng nhận; cân bằng giữa kiểm duyệt chặt (đảm bảo chất lượng) và tốc độ duyệt (không làm nản lòng nhà bán mới).